\section{A product-moment lower bound}
\label{lb:section}

The lower bound is a statement about finite product-moment formulas over
\(\C\).  No positivity of the weights, distinctness of the nodes, or
nonvanishing of their coordinates is required.

\begin{theorem}[Product-moment lower bound]
\label{thm:moment}
Let \(k\geq 1\), and let \(b_i\geq 1\) and \(d_i\geq 2b_i\) be
integers for \(1\leq i\leq k\).  Suppose that
\[
  m_i:\C[z_i]_{\leq d_i}\longrightarrow\C
\]
is a linear functional satisfying
\begin{equation}
  m_i(z_i^e)=0\quad(b_i<e\leq d_i),
  \qquad \beta_i:=m_i(z_i^{b_i})\neq 0.
  \label{lb:truncation}
\end{equation}
Let \(n\geq1\), and let \(z_\ell=(z_{1,\ell},\ldots,z_{k,\ell})\in\C^k\) and
\(w_\ell\in\C\setminus\{0\}\), for \(1\leq\ell\leq n\), satisfy
\begin{equation}
  \sum_{\ell=1}^{n}w_\ell\prod_{i=1}^{k}z_{i,\ell}^{e_i}
  =\prod_{i=1}^{k}m_i(z_i^{e_i})
  \qquad(0\leq e_i\leq d_i).
  \label{lb:moments}
\end{equation}
Set
\begin{equation}
  q_i=b_i+1,\qquad \delta_i=d_i-2b_i,\qquad
  Q=\prod_{i=1}^{k}q_i,\qquad
  N=Q+\sum_{i=1}^{k}\delta_i\frac{Q}{q_i}.
  \label{lb:parameters}
\end{equation}
Then \(n\geq N\).  The nodes may repeat, and any of their coordinates
may be zero.
\end{theorem}

We prove the theorem by constructing \(N\) vectors with an invertible
Gram matrix.  The construction uses a parameter-dependent symmetric
bilinear pairing and indexed evaluation lists.  We first establish
nonsingularity over a rational function field, then choose one common
finite parameter vector.

Introduce independent parameters \(t_1,\ldots,t_k\), and write
\[
  K=\C(t_1,\ldots,t_k),\qquad
  h_{i,\ell}=1+t_i z_{i,\ell},\qquad
  w_\ell(t)=w_\ell\prod_{i=1}^{k}h_{i,\ell}^{d_i}.
\]
Each \(h_{i,\ell}\) is a nonzero polynomial, including when
\(z_{i,\ell}=0\).  On \(K^n\) use the diagonal bilinear pairing
\begin{equation}
  \langle f,g\rangle_t
     =\sum_{\ell=1}^{n}w_\ell(t)f(\ell)g(\ell).
  \label{lb:pairing}
\end{equation}
For \(0\leq\varepsilon_i\leq b_i\), define
\begin{equation}
  u_\varepsilon(\ell)
    =\prod_{i=1}^{k}
       \frac{z_{i,\ell}^{\varepsilon_i}}{h_{i,\ell}^{b_i}},
  \qquad
  \mathcal U=(u_\varepsilon)_{0\leq\varepsilon_i\leq b_i}.
  \label{lb:U}
\end{equation}
For each \(i\) with \(\delta_i>0\), also define
\begin{equation}
  v_{i,s,\varepsilon}(\ell)
    =\frac{z_{i,\ell}^{s}}{h_{i,\ell}^{s}}
       \prod_{j\neq i}
       \frac{z_{j,\ell}^{\varepsilon_j}}{h_{j,\ell}^{b_j}},
  \quad
  \begin{gathered}
    b_i+1\leq s\leq d_i-b_i,\\[-2pt]
    0\leq\varepsilon_j\leq b_j\quad(j\neq i).
  \end{gathered}
  \label{lb:V}
\end{equation}
Denote this list by \(\mathcal V_i\); when \(\delta_i=0\), let
\(\mathcal V_i\) be empty.  Thus \(\mathcal U\) contains \(Q\) vectors
and \(\mathcal V_i\) contains \(\delta_iQ/q_i\) vectors.  Empty products,
including empty Kronecker products of matrices, are interpreted as \(1\).

By linearity, \eqref{lb:moments} applies to any product of univariate
polynomials whose degree in coordinate \(i\) is at most \(d_i\).  We
use it in this form throughout the Gram calculations.

\begin{lemma}
\label{lb:U-gram}
For \(0\leq p,r\leq b_i\), set
\[
  G_i(t_i)_{p,r}
   =m_i\bigl(z_i^{p+r}(1+t_i z_i)^{\delta_i}\bigr).
\]
The Gram matrix of \(\mathcal U\) is
\(\bigotimes_{i=1}^{k}G_i(t_i)\), and
\begin{equation}
  \det G_i(t_i)
    =(-1)^{q_i(q_i-1)/2}\beta_i^{q_i}\neq 0.
  \label{lb:G-det}
\end{equation}
In particular, this Gram matrix is invertible for every finite choice
of the parameters for which the vectors are defined.
\end{lemma}

\begin{proof}
The polynomial defining an entry of \(G_i\) has degree at most
\(2b_i+\delta_i=d_i\).  If \(p+r>b_i\), every monomial in it has
zero \(m_i\)-value by \eqref{lb:truncation}.  If \(p+r=b_i\), only its
lowest term survives, and the entry equals \(\beta_i\).  Reversing the
columns therefore gives an upper triangular matrix with constant
diagonal \(\beta_i\), proving \eqref{lb:G-det}.  Finally,
\eqref{lb:pairing} and \eqref{lb:U} give
\[
  \langle u_\varepsilon,u_\eta\rangle_t
   =\sum_\ell w_\ell\prod_i
      z_{i,\ell}^{\varepsilon_i+\eta_i}h_{i,\ell}^{\delta_i}
   =\prod_iG_i(t_i)_{\varepsilon_i,\eta_i},
\]
which identifies the Gram matrix as the stated Kronecker product.
\end{proof}

\begin{lemma}
\label{lb:orthogonality}
The lists \(\mathcal U,\mathcal V_1,\ldots,\mathcal V_k\) are mutually
orthogonal for \eqref{lb:pairing}.
\end{lemma}

\begin{proof}
In pairing \(v_{i,s,\varepsilon}\) with \(u_\eta\), the factor in
coordinate \(i\), after including the weight, is
\[
   z_i^{s+\eta_i}(1+t_i z_i)^{d_i-s-b_i}.
\]
The exponent of \(1+t_i z_i\) is nonnegative.  Every monomial has
degree at least \(s>b_i\) and at most
\(d_i-b_i+\eta_i\leq d_i\), so its \(m_i\)-value is zero.  Every
other coordinate also has degree at most its prescribed \(d_j\).
Equation~\eqref{lb:moments} proves
\(\langle\mathcal V_i,\mathcal U\rangle_t=0\).

If \(i\neq j\), pair a vector from \(\mathcal V_i\) with one from
\(\mathcal V_j\).  The coordinate-\(i\) factor has the same form,
with \(\eta_i\) supplied by the second vector, and again has zero
moment.  The coordinate-\(j\) factor has nonnegative exponent
\(d_j-s'-b_j\) on \(1+t_jz_j\) and degree at most \(d_j\).
All remaining coordinates have the form used in
Lemma~\ref{lb:U-gram}.  Thus every use of \eqref{lb:moments} stays
inside its degree box, and
\(\langle\mathcal V_i,\mathcal V_j\rangle_t=0\).
\end{proof}

It remains to analyze the Gram matrix within each nonempty
\(\mathcal V_i\).  Put \(r_i=Q/q_i\), and let \(H_i(t)\) be the
\(r_i\)-by-\(r_i\) matrix indexed by the tuples
\(\varepsilon,\eta\) on coordinates other than \(i\), with entries
\begin{equation}
  H_i(t)_{\varepsilon,\eta}
    =\sum_{\ell=1}^{n}w_\ell
       \frac{z_{i,\ell}^{d_i+1}}{1+t_i z_{i,\ell}}
       \prod_{j\neq i}
          z_{j,\ell}^{\varepsilon_j+\eta_j}
          (1+t_jz_{j,\ell})^{\delta_j}.
  \label{lb:H}
\end{equation}

\begin{lemma}
\label{lb:anti-triangular}
Up to the sign of a block-column permutation, the Gram determinant of
\(\mathcal V_i\) equals \(\det(H_i(t))^{\delta_i}\).
\end{lemma}

\begin{proof}
Order the blocks by increasing \(s\).  The \((s,s')\) block has entries
\begin{equation}
  \sum_{\ell=1}^{n}w_\ell
     z_{i,\ell}^{s+s'}h_{i,\ell}^{d_i-s-s'}
     \prod_{j\neq i}
       z_{j,\ell}^{\varepsilon_j+\eta_j}h_{j,\ell}^{\delta_j}.
  \label{lb:V-block}
\end{equation}
If \(s+s'\leq d_i\), this is a polynomial expression in the moment
box, and all its coordinate-\(i\) degrees exceed \(b_i\).  Hence
the entire block is zero.  If \(s+s'=d_i+1\), the block is exactly
\(H_i(t)\).

Write \(s=b_i+1+a\), where \(0\leq a<\delta_i\).  After reversing
the block columns, the column with index \(a'\) corresponds to
\(s'=d_i-b_i-a'\), so
\[
  s+s'=d_i+1+a-a'.
\]
Every block above the diagonal is therefore zero, while every diagonal
block is \(H_i(t)\).  The resulting block lower triangular matrix has
determinant \(\det(H_i(t))^{\delta_i}\), as asserted.
\end{proof}

The next identity proves that \(H_i\) is nonsingular as a rational
matrix.  Its form is chosen to include zero node coordinates.

\begin{lemma}
\label{lb:Laurent}
For integers \(d\geq2b\) and \(b\geq1\), the following identity holds
whenever \(t\neq0\) and \(1+tz\neq0\):
\begin{equation}
  \frac{z^{d+1}}{1+tz}
   =\sum_{p=b}^{d}(-1)^{d-p}
       \frac{z^p}{t^{d+1-p}}
     +(-1)^{d+1-b}
       \frac{z^b}{t^{d+1-b}(1+tz)}.
  \label{lb:Laurent-identity}
\end{equation}
For each fixed \(z\neq0\), the final term has order
\(O(t^{-(d-b+2)})\) at infinity.  For \(z=0\), every term in
\eqref{lb:Laurent-identity} vanishes identically.
\end{lemma}

\begin{proof}
Multiply \eqref{lb:Laurent-identity} by
\((1+tz)t^{d+1-b}\).  Consecutive terms of the finite sum cancel,
leaving exactly the numerator on the left.  If \(z\neq0\), the
additional denominator \(1+tz\) gives one further power of
\(t^{-1}\) at infinity.  If \(z=0\), the factor \(z^b\) vanishes
because \(b\geq1\), as do all the other numerators.  Thus the
identity and the remainder statement require no division by a node
coordinate.
\end{proof}

\begin{lemma}
\label{lb:H-nonsingular}
For every \(i\) with \(\delta_i>0\), \(\det H_i(t)\) is a nonzero
element of \(K\).
\end{lemma}

\begin{proof}
Expand \eqref{lb:H} at \(t_i=\infty\), treating the other parameters
as elements of the coefficient field \(\C(t_j:j\neq i)\).
Apply \eqref{lb:Laurent-identity} separately to each node.  The sum is
finite, and its zero-coordinate terms vanish identically, so the
remainder estimate applies entrywise.  Every term with
\(b_i<p\leq d_i\) has zero coefficient by \eqref{lb:moments} and
\eqref{lb:truncation}.  The coefficient with \(p=b_i\) is
\(\beta_i\bigotimes_{j\neq i}G_j(t_j)\).  Consequently
\begin{equation}
  H_i(t)
   =(-1)^{d_i-b_i}\beta_i\,
       t_i^{-(d_i-b_i+1)}
          \bigotimes_{j\neq i}G_j(t_j)
      +O\bigl(t_i^{-(d_i-b_i+2)}\bigr).
  \label{lb:H-leading}
\end{equation}
This may equivalently be read as a formal Laurent expansion in
\(t_i^{-1}\).  By Lemma~\ref{lb:U-gram}, its leading coefficient
matrix is invertible.  The coefficient of
\(t_i^{-(d_i-b_i+1)r_i}\) in \(\det H_i\) is therefore the nonzero
determinant of that matrix.  In particular, \(\det H_i\) cannot be
the zero rational function.  Each coordinate is treated separately;
no common limiting path for the parameters is required.
\end{proof}

\begin{proof}[Proof of Theorem~\ref{thm:moment}]
Choose one finite parameter vector on which every required block is
nonsingular.  To justify that choice explicitly, put
\[
  D_i(t_i)=\prod_{\ell=1}^{n}(1+t_i z_{i,\ell}).
\]
For \(\delta_i>0\), the entries of \(D_iH_i\) are polynomials in
all the parameters.  Hence
\[
  P_i(t):=D_i(t_i)^{r_i}\det H_i(t)
\]
is a nonzero polynomial by Lemma~\ref{lb:H-nonsingular}.  The product
\[
  \left(\prod_{i=1}^{k}D_i(t_i)\right)
  \left(\prod_{i:\,\delta_i>0}P_i(t)\right)
\]
is nonzero.  Since \(\C\) is infinite, it is nonvanishing at some
finite \(t\in\C^k\).  Fix such a choice.  All the vectors
\eqref{lb:U}--\eqref{lb:V} and the pairing \eqref{lb:pairing} are
then defined.  The Gram matrix of \(\mathcal U\) is invertible by
Lemma~\ref{lb:U-gram}, and each nonempty \(\mathcal V_i\) has an
invertible Gram matrix by Lemmas~\ref{lb:anti-triangular} and
\ref{lb:H-nonsingular}.  Lemma~\ref{lb:orthogonality} makes their
combined Gram matrix block diagonal and invertible.

The combined list thus consists of
\(Q+\sum_i\delta_iQ/q_i=N\) linearly independent vectors in
\(\C^n\), proving \(n\geq N\).  No node was discarded or required
to be distinct or to have nonzero coordinates.  If a moment formula
initially contains zero weights, those terms may first be deleted.
\end{proof}
